\subsection{Contesto e motivazione}

La simulazione numerica accurata ed efficiente di equazioni alle derivate parziali (Partial Differential Equations, PDE) dipendenti dal tempo costituisce un elemento centrale della matematica computazionale e dell'ingegneria scientifica \citep{quarteroni2009the}. In numerosi problemi applicativi, quali il trasporto convettivo, la fluidodinamica e la diffusione di quantità fisiche in presenza di strutture spaziali localizzate, la soluzione può presentare gradienti ripidi, fronti mobili e regioni caratterizzate da scale spaziali molto differenti \citep{donea2003finite}. In tali condizioni, una discretizzazione uniforme e statica può risultare fortemente inefficiente, poiché per risolvere adeguatamente le strutture più piccole è necessario mantenere un'elevata risoluzione anche nelle regioni del dominio in cui la soluzione varia lentamente.

L'adattività della discretizzazione costituisce una delle principali strategie per affrontare questo problema \citep{bangerth2003adaptive}. L'obiettivo è distribuire i gradi di libertà in modo non uniforme, aumentando la risoluzione nelle regioni in cui la soluzione presenta maggiore complessità e riducendola altrove. Nell'ambito dei metodi agli elementi finiti (Finite Element Method, FEM), tale principio può essere perseguito attraverso differenti forme di adattività, tra cui $h$-, $p$- e $r$-refinement.

\subsection{\texorpdfstring{Adattività di griglia e ($r$-refinement)}{Adattività di griglia e(r-refinement)}}

Nel caso dell'$h$-refinement, la risoluzione viene modificata attraverso la suddivisione o il riaccorpamento locale degli elementi, determinando una variazione del passo di discretizzazione e, generalmente, del numero di gradi di libertà \citep{berger1984adaptive}. Sebbene tale approccio sia ampiamente consolidato, nelle simulazioni transitorie caratterizzate da frequenti eventi di adattamento la variazione della connettività può introdurre costi aggiuntivi associati alla gestione della mesh, alla costruzione dei sistemi sparsi e, in implementazioni parallele, al bilanciamento del carico.

Il $p$-refinement modifica invece localmente l'ordine polinomiale delle funzioni di forma mantenendo invariata la geometria degli elementi \citep{szabo2021finite}. Esso può risultare particolarmente efficace per soluzioni sufficientemente regolari, mentre la sua efficacia può ridursi in presenza di strutture non lisce o di forti localizzazioni della soluzione.

L'$r$-refinement, o \textit{moving mesh adaptation}, adotta una strategia differente: il numero di nodi e la connettività topologica vengono mantenuti fissi, mentre le coordinate dei nodi vengono adattate nel tempo per concentrare la risoluzione nelle regioni di maggiore interesse \citep{huang2010adaptive,budd2009adaptivity,hyman1988moving}. Questa caratteristica rende l'$r$-refinement particolarmente interessante per problemi transitori, poiché consente di modificare la distribuzione spaziale dei gradi di libertà senza alterare la dimensione del sistema algebrico associato alla discretizzazione.

I metodi classici di moving mesh determinano generalmente la nuova configurazione della mesh attraverso formulazioni variazionali, principi di equidistribuzione o equazioni ausiliarie per il moto della griglia, tra cui le Moving Mesh Partial Differential Equations (MMPDE) \citep{huang2001practical,huang2010adaptive,budd2009adaptivity}. Sebbene tali formulazioni siano consolidate e possano produrre mesh di elevata qualità, il calcolo della configurazione adattata costituisce un problema aggiuntivo rispetto alla PDE fisica. Quando l'adattamento deve essere eseguito frequentemente, il costo della procedura di mesh movement può quindi diventare una componente significativa del tempo complessivo di simulazione \citep{song2022m2n,rowbottom2024g}.

\subsection{Metodi di mesh movement basati su apprendimento automatico}

Lo sviluppo dello Scientific Machine Learning e del Geometric Deep Learning ha recentemente motivato l'impiego di modelli neurali per apprendere operatori definiti direttamente su mesh non strutturate e grafi \citep{bronstein2021geometric,battaglia2018relational}. In questo contesto, una mesh agli elementi finiti può essere rappresentata naturalmente mediante un grafo i cui nodi corrispondono ai gradi di libertà geometrici e i cui archi codificano le relazioni di adiacenza.

Un primo filone di lavori ha mostrato che reti neurali possono essere utilizzate per sostituire, o accelerare, le procedure classiche di mesh movement. M2N, ad esempio, introduce un framework basato su apprendimento automatico per la rilocazione dei nodi e considera esplicitamente problematiche quali il mesh tangling, la consistenza al bordo e la generalizzazione a differenti risoluzioni; il metodo viene validato sia su problemi stazionari sia su problemi tempo-dipendenti \citep{song2022m2n}. Successivamente, DMM e il relativo framework MM-PDE hanno proposto un mesh mover basato su una formulazione neurale della mappa di adattamento, mantenendo fissi i nodi esistenti e integrando il movimento della mesh all'interno di un solutore per PDE dinamiche \citep{hu2024better}.

Un ulteriore sviluppo è rappresentato da UM2N, che impiega un'architettura basata su Graph Transformer e Graph Attention Network per apprendere il movimento della mesh e ne studia l'applicabilità a differenti PDE, geometrie e risoluzioni, includendo problemi di advezione e Navier--Stokes \citep{zhang2024towards}. Più recentemente, G-Adaptivity ha formulato direttamente il problema della rilocazione dei nodi nel contesto FEM come un problema di apprendimento mediante GNN, ottimizzando la posizione dei nodi in funzione dell'errore della soluzione agli elementi finiti e mirando a ridurre il costo delle procedure classiche di adattamento \citep{rowbottom2024g}. Un lavoro recente ha inoltre studiato l'impiego di reti neurali per $r$-adaptive FEM in problemi ellittici autoaggiunti parametrizzati, evidenziando tuttavia un'impostazione specifica per problemi coercivi e per mesh aventi struttura tensoriale \citep{aballay2026radaptive}.

Nel loro insieme, questi lavori mostrano che il machine learning può ridurre significativamente il costo della generazione di mesh adattive e che le GNN costituiscono una scelta naturale per operare direttamente sulla struttura della discretizzazione. Essi mostrano inoltre che la gestione della qualità geometrica della mesh rappresenta una componente essenziale dei metodi neurali di mesh movement. In particolare, il problema della possibile inversione o degenerazione degli elementi è stato affrontato mediante differenti strategie, tra cui trasformazioni neurali strutturalmente invertibili, termini di perdita geometrici e architetture progettate per limitare il mesh tangling \citep{song2022m2n,zhang2024towards}.

Questi risultati evidenziano tuttavia una distinzione importante. La predizione di una nuova configurazione della mesh non specifica, di per sé, come l'evoluzione della soluzione numerica debba essere trattata quando la geometria cambia nel tempo. In particolare, occorre distinguere tra una sequenza di rilocazioni discrete della mesh, che richiede il trasferimento dello stato numerico tra spazi FEM differenti, e una formulazione in cui il movimento della mesh entra direttamente nell'integrazione temporale della PDE.

\subsection{Remapping, mesh mobile e formulazione ALE}

Quando una procedura di adattamento produce una nuova configurazione $\mathcal{T}^{n+1}$ a partire da $\mathcal{T}^{n}$ mediante una rilocazione discreta dei nodi, lo stato numerico rappresentato sulla mesh precedente deve essere trasferito al nuovo spazio discreto. Indicando con $u^n_{\mathcal{T}^n}$ la soluzione numerica sulla configurazione precedente, il trasferimento assume schematicamente la forma
\[
u^n_{\mathcal{T}^n}
\longrightarrow
u^n_{\mathcal{T}^{n+1}}.
\]
Tale operazione di \textit{state remapping} introduce un ulteriore operatore numerico nella pipeline e, nel caso di simulazioni con numerosi eventi di adattamento, il relativo errore può accumularsi nel tempo. La scelta dell'operatore di trasferimento è pertanto parte integrante della qualità numerica complessiva del metodo.

Una possibile strategia consiste nell'utilizzare un'interpolazione nodale o elementare; un'alternativa più accurata è rappresentata dalla proiezione variazionale, ad esempio mediante una proiezione in norma $L^2$, che consente di trasferire lo stato imponendo una condizione di ortogonalità rispetto al nuovo spazio FEM. Tuttavia, una proiezione di questo tipo richiede la costruzione e la soluzione di un sistema associato alla nuova matrice di massa, introducendo un costo algebrico non trascurabile. La questione non consiste quindi soltanto nell'ottenere una mesh adattata, ma nel determinare come evolvere la soluzione numerica durante il movimento della discretizzazione senza introdurre un costo di trasferimento eccessivo.

Una formulazione alternativa è costituita dal paradigma Arbitrary Lagrangian--Eulerian (ALE), nel quale la mesh viene descritta mediante una mappa tempo-dipendente
\[
\mathbf{x}=\mathbf{x}(\boldsymbol{\xi},t)
\]
e il suo moto è caratterizzato dalla velocità di griglia
\[
\mathbf{w}=\frac{\partial \mathbf{x}}{\partial t}.
\]
La velocità della mesh modifica direttamente i termini convettivi della formulazione differenziale, evitando di trattare il cambiamento della configurazione come una successione di trasferimenti indipendenti dello stato \citep{hirt1974an,donea2004arbitrary}.

L'impiego di una mesh mobile introduce tuttavia una condizione di consistenza geometrica fondamentale. La discreta evoluzione della geometria e l'integrazione temporale della PDE devono essere compatibili con la Geometric Conservation Law (GCL), affinché un campo uniforme possa essere preservato durante il moto della griglia. Nel contesto FEM-ALE, la D-GCL (Discrete Geometric Conservation Law) fornisce quindi una condizione discreta per la corretta rappresentazione dell'evoluzione geometrica e per la preservazione di stati uniformi sulla mesh deformante \citep{farhat2001the,formaggia2004stability}.

È importante distinguere tale proprietà dalla conservazione delle grandezze fisiche della PDE: la verifica della D-GCL e della preservazione del campo costante costituisce una verifica della consistenza geometrica della discretizzazione mobile, e non una dimostrazione generale di conservazione fisica di massa, quantità di moto o energia.

\subsection{Problema aperto e obiettivo del lavoro}

La letteratura esistente ha quindi sviluppato, da un lato, metodi classici e neurali per la generazione e la rilocazione di mesh adattive e, dall'altro, formulazioni FEM-ALE mature per la soluzione di PDE su griglie deformanti \citep{budd2009adaptivity,song2022m2n,hu2024better,zhang2024towards,rowbottom2024g}. Questi filoni affrontano tuttavia problemi differenti: i metodi di mesh movement learning-based si concentrano principalmente sulla predizione della configurazione geometrica adattata e sul costo della procedura di mesh adaptation, mentre la formulazione ALE riguarda il modo in cui il moto della geometria viene incorporato nel solutore differenziale.

Nella letteratura da noi esaminata non abbiamo identificato una pipeline FEM transiente che combini esplicitamente tre componenti in un unico processo numerico: 
(i) la predizione neurale del movimento nodale mediante una GNN;
(ii) una verifica geometrica deterministica e indipendente dalla rete che certifichi l'ammissibilità dello spostamento rispetto alla positività del Jacobiano elementare;
e (iii) l'integrazione del moto predetto come velocità di griglia in una formulazione ALE continua verificata rispetto alla D-GCL, evitando il ricorso al remapping esplicito dello stato ad ogni evento di adattamento.

L'obiettivo del presente lavoro è quindi studiare questa integrazione in modo sistematico, non assumendo a priori che ciascuna componente sia sufficiente isolatamente. Il framework proposto viene pertanto analizzato attraverso una sequenza di verifiche disaccoppiate, finalizzate a distinguere l'errore associato alla predizione della mesh, quello introdotto dal trasferimento dello stato e quello associato alla discretizzazione su mesh mobile.

\subsection{Contributi del lavoro}

In questo articolo presentiamo \textit{Graph Dynamic Meshing} (GDM), un framework ibrido neurale e numerico per l'$r$-adattività di mesh bidimensionali non strutturate a topologia fissa. Il framework combina un modello GNN per la predizione del movimento della mesh, un controllo geometrico deterministico post-inferenza e una formulazione ALE per l'evoluzione della soluzione sulla griglia mobile.

I principali contributi del lavoro sono:

\begin{enumerate}

    \item \textbf{Predizione neurale del movimento nodale:}
    sviluppiamo una rete neurale su grafi che opera direttamente sulla struttura della mesh e predice gli spostamenti nodali sulla base di informazioni geometriche e fisiche disponibili durante la simulazione. La procedura consente di sostituire, in fase di inferenza, la risoluzione iterativa dell'oracolo variazionale con una singola valutazione del modello. Lo scaling computazionale dell'inferenza viene inoltre analizzato sperimentalmente al variare della dimensione della mesh.

    \item \textbf{Filtro geometrico deterministico post-inferenza:}
    introduciamo un \textit{Deterministic Safety Filter} separato dalla rete neurale. Per elementi triangolari lineari e per una traiettoria lineare dello spostamento, il determinante Jacobiano, equivalente alla variazione dell'area orientata, è una funzione quadratica del fattore di scala $\alpha$. Le sue radici possono quindi essere analizzate in forma chiusa per determinare il massimo fattore globale di contrazione ammissibile, garantendo
    \[
    \min_{e} J_e(\alpha) \geq \epsilon>0
    \]
    lungo la traiettoria considerata. La procedura fornisce quindi una verifica geometrica deterministica dell'output neurale, indipendente dalla funzione di perdita utilizzata durante l'addestramento.

    \item \textbf{Accoppiamento con una formulazione ALE continua:}
    trasformiamo lo spostamento nodale predetto in una velocità di griglia e integriamo il relativo moto direttamente nel solutore FEM mediante una formulazione ALE. In questo modo il cambiamento della geometria non richiede, nella pipeline ALE principale, un remapping esplicito dello stato numerico. La consistenza geometrica della formulazione viene verificata mediante test dedicati di D-GCL e di preservazione di un campo uniforme.

    \item \textbf{Validazione numerica disaccoppiata:}
    conduciamo una campagna di esperimenti su una PDE di convezione--diffusione transiente con Manufactured Solution, separando sistematicamente la verifica del mesh mover, del filtro geometrico, degli operatori di remapping e della formulazione ALE. Vengono inoltre analizzati gli ordini di convergenza spaziale e temporale, la sensibilità alla frequenza di adattamento e il costo computazionale delle differenti componenti.

\end{enumerate}

\subsection{Organizzazione del paper}

Il resto dell'articolo è organizzato come segue.

La Sezione~\ref{sec:formulation} presenta la formulazione matematica della PDE di riferimento, la discretizzazione mediante elementi finiti $P_1$ con time-stepping Backward Euler e la definizione del funzionale variazionale utilizzato come oracolo per la generazione dei target di adattamento.

La Sezione~\ref{sec:gdm_framework} descrive il framework GDM, includendo la rappresentazione della mesh come grafo, l'architettura della rete \textit{MeshMoverGNN} e la derivazione del Deterministic Safety Filter.

La Sezione~\ref{sec:ale_dgcl} introduce la formulazione ALE utilizzata per l'evoluzione su mesh mobile e presenta la relativa verifica della D-GCL. La sezione discute inoltre il ruolo del remapping nei metodi a rilocazione discreta e la sua sostituzione, nella pipeline ALE, mediante l'evoluzione continua della geometria.

La Sezione~\ref{sec:diagnostics} riporta le campagne di verifica disaccoppiata dei singoli sottosistemi, comprendendo i risultati relativi al mesh mover, al filtro geometrico, al remapping e alle proprietà di convergenza della discretizzazione.

La Sezione~\ref{sec:dynamic_benchmark} presenta il benchmark dinamico finale e il confronto tra le differenti strategie di adattamento e di evoluzione della mesh, analizzando accuratezza della soluzione, errore integrato nel tempo e costo computazionale.

La Sezione~\ref{sec:discussion} discute i risultati, i compromessi tra accuratezza e costo computazionale e le principali limitazioni dell'attuale formulazione.

Infine, la Sezione~\ref{sec:conclusions} riassume i risultati ottenuti e delinea le prospettive di sviluppo verso problemi più complessi e PDE non lineari.
